\section{Approach: Claim-Survival Analysis}\label{sec:approach}

\textbf{Construct and decomposition.} An explanation $E$ of function $f$ is written at anchor $c_0$. The prompt requires each claim to be a self-contained, source-checkable fact naming $f$, with no added or merged facts, original wording where possible, and a 12-claim cap. Each receives one type: \emph{structural} (signature, parameters, calls), \emph{data-flow}, \emph{behavioural} (conditions, exceptions, side effects), or \emph{intent}. These instructions do not guarantee atomicity or coverage. A \emph{touch} $v_k$ is a later mainline commit changing $f$'s code tokens (docstrings, comments and layout removed). Git file renames are followed; fallback matching requires a unique new same-name candidate in the same file with text similarity $\ge0.7$. This is an operational tracking rule, not validated identity resolution: false matches may create spurious follow-up, while failed matches censor tracks and can be informative. A judge $J(c,f@v_k)\in\{\textsc{sup},\textsc{con},\textsc{unr}\}$ uses docstring-free code, the diff and tree-sitter~\cite{treesitter} facts (parameters, calls, returns, raises).

\textbf{Cohort, event, censoring.} Only anchor-\textsc{sup} claims enter; anchor contradictions are initial detected errors. The event is the \emph{first} \textsc{con}, so later support cannot undo it. Follow-up is censored at $K_{\max}$, the pinned head, or unmatched removal/rename. Kaplan--Meier~\cite{kaplan1958} estimates $S(k)=\Pr[T>k]$ for this detection time, conditional on anchor support. We report $\hat S(K)$ and the time to 10\% detections; the half-life is not reached. \emph{Unresolved verdicts:} primary analysis keeps \textsc{unr} claims at risk without an event; sensitivity scenarios instead censor immediately before the first \textsc{unr}, or count it as an event. Neither scenario bounds semantic survival; unresolved is not supported. Formal error adjustment needs judge sensitivity/specificity that this archive cannot identify. Follow-up stops at the first detection, precluding analysis of later recovery.

\textbf{Explaining detections.} To explore which edits accompany detections, we fit a discrete-time complementary log-log model with a separate baseline for each touch index. Predictors are claim type, $\log(1{+}\text{lines changed})$ in the function and file, and change flags derived from parser facts. A hazard ratio (HR) describes the relative detection rate among claims still at risk. Intervals and Wald tests account for claims sharing a function. Cox models with changing covariates~\cite{cox1972,andersen1982,davidson2019lifelines} and proportional-hazards checks are retained as robustness analyses in the artifact. 

\textbf{Placebo sensitivity.} Rejudging unchanged code gives a scenario with constant false-contradiction rate $\epsilon$: $\max(0,(h_k-\epsilon)/(1-\epsilon))$. Error adjustment would require placebo rates to transfer to real-diff prompts, independent repeated errors conditional on risk, and no missed contradictions. Contexts differ and judges disagree both ways; these assumptions are unvalidated. The curve is a sensitivity scenario, not a calibrated correction.

\textbf{Exploratory comparators.} A whole-explanation JIT check by the same model uses three prompt versions (binary v1 and v2, graded 0--10 v3); a non-LLM \emph{identifier check} flags a claim once an identifier it names, present at $c_0$, is missing from the code. We also decompose and judge the function's anchor docstring in the same shuffled list as the LLM claims. These are implementation diagnostics, not published-baseline comparisons or causal authorship tests.
